\medskip




S'il choisit de prendre les transports en commun un matin, il reprend les transports en
commun le lendemain avec une probabilité égale à 0,8.

S'il utilise son vélo un matin, il reprend son vélo le lendemain avec une probabilité égale
à 0,4.

\medskip

Pour tout entier naturel $n$ non nul, on note:
\begin{itemize}
\item  $T_n$ l'évènement \og Monsieur Durand utilise les transports en commun le $n$-ième jour\fg
\item  $V_n$ l'évènement \og  Monsieur Durand utilise son vélo le $n$-ième jour\fg
\item On note $p_n$ la probabilité de l'évènement $T_n$,
\end{itemize}

Le premier matin, il décide d'utiliser les transports en commun. Ainsi, la probabilité de
l'évènement $T_1$ est $p_1 = 1$.

\begin{enumerate}
\item 
2\ieme{} et 3\ieme{} jours,

\begin{center}
\psset{nodesepA=0pt,nodesepB=3pt,treesep=0.75,levelsep=3cm,labelsep=0.1pt}
\pstree[treemode=R]{\TR{$T_1$~}}
{\pstree{\TR{$T_2$~}\taput{$0,8$}}
	{
	\TR{$T_3$~}\taput{$0,8$}
	\TR{$V_3$~}\tbput{$0,2$}
	}
\pstree{\TR{$V_2$~}\tbput{$0,2$}}
	{\TR{$T_3$~}\taput{$0,6$}
	\TR{$V_3$~}\tbput{$0,4$}
	}
}
\end{center}

\item Calculer $p_3$

D'après la loi des probabilités totales : $p_3 = p\left(T_2 \cap T_3 \right)  + p\left(V_2 \cap T_3\right)$.

$p\left(T_2 \cap T_3 \right) = p\left(T_2\right) \times p_{T_2}\left(T_3\right)= 0,8 \times 0,8 = 0,64$. De même :

$p\left(V_2 \cap T_3\right) = p\left(V_2\right) \times p_{V_2}\left(T_3\right)= 0,2 \times 0,6 = 0,12$.

Donc $p_3 = 0,64 + 0,12 = 0,76$.
\item  Le 3\ieme{} jour, M. Durand utilise son vélo.

Calculer la probabilité qu'il ait pris les transports en commun la veille.

On calcule $p_{V_3}\left(T_2\right) = \dfrac{p\left(V_3 \cap T_2\right)}{p\left(V_3\right)} =
\dfrac{p\left(T_2 \cap V_3\right)}{p\left(V_3\right)}$.

Or $p\left(V_3\right) = 1 - p_3 = 1 - 0,76 = 0,24$ et 

$p\left(T_2 \cap V_3\right) = 0,8 \times 0,2 = 0,16$, d'où 

$p_{V_3}\left(T_2\right) = \dfrac{0,16}{0,24} = \dfrac{16}{24} = \dfrac{2}{3}$.

\item Recopier et compléter l'arbre pondéré ci-dessous représentant la situation pour les
$n$-ième et $(n + 1)$-ième jours.

\begin{center}
\psset{nodesepA=0pt,nodesepB=3pt,treesep=0.75,levelsep=3cm,labelsep=0.1pt}
\pstree[treemode=R]{\TR{}}
{\pstree{\TR{$T_n$~}\taput{$p_n$}}
	{
	\TR{$T_{n+1}$~}\taput{$0,8$}
	\TR{$V_{n+1}$~}\tbput{$0,2$}
	}
\pstree{\TR{$V_n$~}\tbput{$1 - p_n$}}
	{\TR{$T_{n+1}$~}\taput{$0,6$}
	\TR{$V_{n+1}$~}\tbput{$0,4$}
	}
}
\end{center}

\item  Montrer que, pour tout entier naturel $n$ non nul, $p_{n+1} = 0,2p_n + 0,6$.

On a $p_{n+1} = p\left(T_{n+1}\right) = p\left(T_n \cap T_{n+1} \right)  + p\left(V_n \cap T_{n+1}\right) = 0,8p_n + 0,6\left(1 - p_n\right) = 0,8p_n + 0,6 - 0,6p_n = 0,2p_n + 0,6$.
\item  Montrer par récurrence que, pour tout entier naturel $n$ non nul, on a
\[p_n = 0,75 + 0,25 \times 0,2^{n-1}.\]

\emph{Initialisation} : pour $n = 1$, on a $p_1 = 0,75 + 0,25 \times 0,2^{1-1} = 0,75 + 0,25 \times 0,2^0 = 0,72 + 0,25 = 1$ : la relation est vraie au rang 1.

\emph{Hérédité} : soit $n \geqslant 1$ tel que $p_{n+1} = 0,75 + 0,25 \times 0,2^{n-1}$.

Alors d'après la question 5. $p_{n+1} = 0,2p_n + 0,6$, donc $p_{n+2} = 0,2p_{n+1} + 0,6$, d'où en utilisant l'hypothèse de récurrence :

$p_{n+2} = 0,2\left(0,75 + 0,25 \times 0,2^{n-1}\right) + 0,6  = 0,15 + 0,25\times 0,2^{n} + 0,6 = 0,75 + \times 0,2^{n} + 0,6$ : l'égalité est vraie au rang $n + 1$.

\emph{Conclusion} : la relation est vraie au rang et si elle vrai au rang $n$ au moins égal à 1 elle l'est aussi au rang $n + 1$ : d'après le principe de récurrence :

Pour tout entier naturel $n$ non nul, $p_n = 0,75 + 0,25 \times 0,2^{n-1}$.
\item Déterminer la limite de la suite $(p_n)$ et interpréter le résultat dans le contexte de
l'exercice.

Comme $- 1 < 0,2 < 1$, on sait que $\displaystyle\lim_{n \to + \infty}  0,2^{n-1} = 0$, donc $\displaystyle\lim_{n \to + \infty} 0,25\times  0,2^{n-1} = 0$ et par somme de limites $\displaystyle\lim_{n \to + \infty} p_n = 0,75 = \dfrac34$.

Au bout d'un certain nombre de jours Monsieur Durand prendra les transports en commun 3 jours sur 4.
\end{enumerate}

\bigskip
